\documentclass{beamer}

%
% Choose how your presentation looks.
%
% For more themes, color themes and font themes, see:
% http://deic.uab.es/~iblanes/beamer_gallery/index_by_theme.html
%
\mode<presentation>
{
  \usetheme{default}      % or try Darmstadt, Madrid, Warsaw, ...
  \usecolortheme{default} % or try albatross, beaver, crane, ...
  \usefonttheme{default}  % or try serif, structurebold, ...
  \setbeamertemplate{navigation symbols}{}
  \setbeamertemplate{caption}[numbered]
  \setbeamertemplate{footline}[page number]
  \setbeamercolor{frametitle}{fg=white}
  \setbeamercolor{footline}{fg=black}
} 

\usepackage[english]{babel}
\usepackage[utf8x]{inputenc}
\usepackage{tikz}
\usepackage{listings}
\usepackage{courier}
\usepackage{minted}

\xdefinecolor{darkblue}{rgb}{0.1,0.1,0.7}
\xdefinecolor{darkgreen}{rgb}{0,0.5,0}
\xdefinecolor{darkorange}{rgb}{0.8,0.5,0}
\xdefinecolor{darkred}{rgb}{0.7,0,0}
\xdefinecolor{dianablue}{rgb}{0.18,0.24,0.31}
\definecolor{commentgreen}{rgb}{0,0.6,0}
\definecolor{stringmauve}{rgb}{0.58,0,0.82}

\lstset{ %
  backgroundcolor=\color{white},      % choose the background color
  basicstyle=\ttfamily\small,         % size of fonts used for the code
  breaklines=true,                    % automatic line breaking only at whitespace
  captionpos=b,                       % sets the caption-position to bottom
  commentstyle=\color{commentgreen},  % comment style
  escapeinside={\%*}{*)},             % if you want to add LaTeX within your code
  keywordstyle=\color{blue},          % keyword style
  stringstyle=\color{stringmauve},    % string literal style
  showstringspaces=false,
  showlines=true
}

\lstdefinelanguage{scala}{
  morekeywords={abstract,case,catch,class,def,%
    do,else,extends,false,final,finally,%
    for,if,implicit,import,match,mixin,%
    new,null,object,override,package,%
    private,protected,requires,return,sealed,%
    super,this,throw,trait,true,try,%
    type,val,var,while,with,yield},
  otherkeywords={=>,<-,<\%,<:,>:,\#,@},
  sensitive=true,
  morecomment=[l]{//},
  morecomment=[n]{/*}{*/},
  morestring=[b]",
  morestring=[b]',
  morestring=[b]"""
}

\title[2016-12-12-femtocode-at-diana]{Femtocode: development of a query \\ language for HEP}
\author{Jim Pivarski}
\institute{Princeton -- DIANA}
\date{December 12, 2016}

\begin{document}

\logo{\pgfputat{\pgfxy(0.11, 8)}{\pgfbox[right,base]{\tikz{\filldraw[fill=dianablue, draw=none] (0 cm, 0 cm) rectangle (50 cm, 1 cm);}}}\pgfputat{\pgfxy(0.11, -0.6)}{\pgfbox[right,base]{\tikz{\filldraw[fill=dianablue, draw=none] (0 cm, 0 cm) rectangle (50 cm, 1 cm);}\includegraphics[height=0.99 cm]{diana-hep-logo.png}\tikz{\filldraw[fill=dianablue, draw=none] (0 cm, 0 cm) rectangle (4.9 cm, 1 cm);}}}}

\begin{frame}
  \titlepage
\end{frame}

\logo{\pgfputat{\pgfxy(0.11, 8)}{\pgfbox[right,base]{\tikz{\filldraw[fill=dianablue, draw=none] (0 cm, 0 cm) rectangle (50 cm, 1 cm);}\includegraphics[height=1 cm]{diana-hep-logo.png}}}}

% Uncomment these lines for an automatically generated outline.
%\begin{frame}{Outline}
%  \tableofcontents
%\end{frame}

\begin{frame}{Motivations}
\textcolor{darkblue}{\underline{Specific motivation}: to accelerate data pulls, both}

\begin{center}
\begin{columns}
\column{0.1\linewidth}

\column{0.28\linewidth}
\centering \textcolor{darkblue}{for the humans}

{\it shorter time from concept to code}

\column{0.05\linewidth}
\centering \textcolor{darkblue}{and}

\column{0.28\linewidth}
\centering \textcolor{darkblue}{for the computers.}

{\it shorter time from code to results}

\column{0.1\linewidth}
\end{columns}
\end{center}

\vspace{0.2 cm}
\textcolor{gray}{(By ``data pull,'' I mean extract, transform, and filter data from a collaboration's central dataset (AOD) and download it for further processing. Right now, we're writing C++ modules.)}

\vspace{0.5 cm}
\begin{uncoverenv}<2->
\textcolor{darkblue}{``\underline{Chip on my shoulder}''}

\vspace{0.2 cm}
To show that a very high-level, abstract view of computation can nevertheless be fast.
\end{uncoverenv}
\end{frame}

\begin{frame}[fragile]{High-level and speed are not antithetical}
\vspace{0.5 cm}
\begin{columns}
\column{0.55\linewidth}
Code like

\small
\begin{minted}{scala}
bestmuon =
  muons.filter(m => m.iso > 10)
       .maxBy(m => m.pt)
\end{minted}
\vspace{0.4 cm}

\column{0.35\linewidth}
\\ does {\it not} need to create function objects or muon objects at runtime!

\vspace{0.2 cm}
It need not be ``taken literally.''
\end{columns}

\vspace{0.5 cm}
\textcolor{darkblue}{Another possible execution plan:}
\begin{enumerate}
\item Start with all {\tt\small muon.iso} values in one array, all {\tt\small muon.pt} values in another array, and a ``repetition level'' to specify where events begin and end.
\item Apply the contents of the filter function to make a mask.
\item Use the mask and repetition level to compact the {\tt\small muon.pt} into zero or one results per event.
\end{enumerate}
\end{frame}

\begin{frame}{How fast {\it should} it be?}
\vspace{0.3 cm}
\textcolor{darkorange}{\bf Suppose we want to plot all CMS $t\bar{t}$ data ever collected.}
\begin{itemize}
\item About 100~fb$^{-1}$ (\href{http://cms-service-lumi.web.cern.ch/cms-service-lumi/publicplots/int_lumi_cumulative_pp_1.png}{\textcolor{blue}{plot}}) $\times$ 1~nb (\href{https://atlas.web.cern.ch/Atlas/GROUPS/PHYSICS/CombinedSummaryPlots/SM/ATLAS_p_SMSummary_SqrtS_Zoom/ATLAS_p_SMSummary_SqrtS_Zoom.png}{\textcolor{blue}{plot}}) is 100 million events.
\item Derivation of a plotted quantity could involve $\mathcal{O}(100)$ double-precision values per event.
\item That's only 80~GB of data that needs to be evaluated; anything else is overhead.
\end{itemize}

\vspace{0.5 cm}
\textcolor{darkblue}{{\it \underline{Single}} computer, {\it \underline{single}} thread:}

\vspace{0.1 cm}
\begin{description}
\item[load from disk:] 21 minutes (Numpy), 22 minutes (ROOT)
\item[CPU execution:] 70 seconds (no optimization), 13 seconds ({\tt -O3})
\item[GPU execution:] 12 seconds to copy to GPU, 5 seconds to compute
\end{description}

\vspace{0.1 cm}
\textcolor{gray}{(Probably close to ``theoretical optimum.'' Distributed processing and multi-user cache can help.)}
\end{frame}

\begin{frame}{Role of language}
\vspace{0.5 cm}
Directly filling histograms from the collaboration-wide repository, eliminating private skims, would fundamentally improve analysis workflows. Analysis could become ``interactive.''

\vspace{0.5 cm}
\textcolor{darkblue}{But\ldots\ how can a language help?}

\begin{uncoverenv}<2->
\vspace{0.3 cm}
\begin{itemize}
\item Writing vectorized, one-array-at-a-time algorithms is unnatural, distracting from analysis.
\item Physicists {\it should} be thinking about one event at a time and muons as objects.
\end{itemize}

\vspace{0.3 cm}
\textcolor{darkblue}{Transforming human concepts into physical execution is exactly what programming languages are for.}
\end{uncoverenv}
\end{frame}

\begin{frame}{What's wrong with the languages we have?}
\vspace{0.25 cm}
As I said in \href{https://indico.cern.ch/event/544650/}{\textcolor{blue}{my June 20 talk}}, C++ and Python are too expressive to permit these kinds of optimizations.

\vspace{0.5 cm}
Automatic vectorization is (in general) an unsolved problem in computer science. But building a vectorizable language is relatively easy: they're just less capable. Typically, query languages like SQL are vectorized.

\vspace{0.5 cm}
Business intelligence with SQL is typically ``interactive,'' and big data projects are pushing $\mathcal{O}(\mbox{second})$ response times to petabyte scales: Ibis, Impala, Kudu, Drill, \ldots\ others? (\href{http://research.google.com/pubs/pub36632.html}{\textcolor{blue}{Google paper}}).

\vspace{0.5 cm}
\textcolor{darkorange}{\bf However, SQL and its relatives are not sufficient for us.}
\end{frame}

\begin{frame}[fragile]{Nested query in C++}
\vspace{0.5 cm}
{\bf Example query:}
\begin{center}
\begin{minipage}{0.95\linewidth}
\textcolor{darkblue}{``Momentum of the track with $|\eta|$ $<$ 2.4 that has the most hits.''}
\end{minipage}
\end{center}
\small
\begin{minted}{c++}
Track *best = NULL;
for (int i = 0;  i < tracks.size();  i++) {
  if (fabs(tracks[i]->eta) < 2.4)
    if (best == NULL ||
        tracks[i]->hits.size() > best->hits.size())
      best = tracks[i];
}
if (best != NULL)
  return best->pt;
else
  return 0.0;
\end{minted}
\end{frame}

\begin{frame}[fragile]{Nested query in SQL}
\vspace{0.5 cm}
{\bf Example query:}
\begin{center}
\begin{minipage}{0.95\linewidth}
\textcolor{darkblue}{``Momentum of the track with $|\eta|$ $<$ 2.4 that has the most hits.''}
\end{minipage}
\end{center}
\small
\begin{minted}{sql}
WITH hit_stats AS (
  SELECT hit.track_id, COUNT(*) AS hit_count FROM hit
    GROUP BY hit.track_id),
 track_sorted AS (
    SELECT track.*, 
    ROW_NUMBER() OVER (
     PARTITION BY track.event_id
     ORDER BY hit_stats.hit_count DESC)
  track_ordinal FROM track INNER JOIN hit_stats
    ON hit_stats.track_id = track.id
    WHERE ABS(track.eta) < 2.4)
 SELECT * FROM event INNER JOIN track_sorted
   ON track_sorted.event_id = event.id
WHERE
  track_sorted.track_ordinal = 1
\end{minted}
\end{frame}

\begin{frame}[fragile]{Nested query in Femtocode}
\vspace{0.5 cm}
{\bf Example query:}
\begin{center}
\begin{minipage}{0.95\linewidth}
\textcolor{darkblue}{``Momentum of the track with $|\eta|$ $<$ 2.4 that has the most hits.''}
\end{minipage}
\end{center}
\small
\begin{onlyenv}<1>
\begin{minted}{scala}
event.tracks
     .filter(t => abs(t.eta) < 2.4)
     .maxBy(t => t.hits.size)
     .map(t => t.pt)
     .impute(0.0)
\end{minted}
\end{onlyenv}
\begin{onlyenv}<2>
\begin{minted}{bash}
event.tracks
     .filter(abs($1.eta) < 2.4)
     .maxBy($1.hits.size)
     .map($1.pt)
     .impute(0.0)
\end{minted}
\end{onlyenv}
\end{frame}

\begin{frame}[fragile]{Femtocode: tiny, declarative code snippets}
\vspace{0.4 cm}
Combine the {\it implementation flexibility} of declarative languages like SQL with the {\it expressiveness} of a functional language for dealing with nested structure.

\vspace{0.5 cm}
{\bf Example use in Python:}

\small
\begin{minted}{python}
h = db.dataset("ttbar-MC")
  .withColumn(varName = "<Femtocode goes here>")
  .filter("<Femtocode using varName>")
  .flatMap("<Femtocode changing nesting level>")
  .Label(hist1 = Bin(100, -5.0, 5.0, "<Femtocode>"),
         hist2 = Bin(20, 0.0, 100.0, "<Femtocode>"),
         hist3 = Bin(314, -pi, pi, "<Femtocode>"))
\end{minted}

\normalsize
followed by Python analysis on {\small\tt h[\textcolor{darkred}{"hist1"}]}, {\small\tt h[\textcolor{darkred}{"hist2"}]}\ldots

\vspace{0.1 cm}
\begin{uncoverenv}<2->
Everything in quotes is the statically checked and optimized big data pull; everything outside is flexible Python.

\vspace{0.1 cm}
Histogrammar ({\tt\small Label} and {\tt\small Bin}) is a necessary component.
\end{uncoverenv}
\end{frame}

\begin{frame}[fragile]{Working examples}
\vspace{0.3 cm}
\begin{lstlisting}[language=python]
>>> from femtocode.parser import parse
>>> print ast.dump(parse("""
... tracks.filter(abs($1.eta) < 2.4)
...       .maxBy($1.hits.size)
...       .map($1.pt)
...       .impute(0.0)
... """))
\end{lstlisting}
\begin{lstlisting}[language=python,basicstyle=\ttfamily\scriptsize]
Suite(assignments=[], expression=FcnCall(function=Attribute(value=FcnCall(function=Attribute(value=FcnCall(function=Attribute(value=FcnCall(function=Attribute(value=Name(id='tracks', ctx=Load()), attr='filter', ctx=Load()), positional=[Compare(left=FcnCall(function=Name(id='abs', ctx=Load()), positional=[Attribute(value=AtArg(num=1), attr='eta', ctx=Load())], names=[], named=[]), ops=[Lt()], comparators=[Num(n=2.4)])], names=[], named=[]), attr='maxBy', ctx=Load()), positional=[Attribute(value=Attribute(value=AtArg(num=1), attr='hits', ctx=Load()), attr='size', ctx=Load())], names=[], named=[]), attr='map', ctx=Load()), positional=[Attribute(value=AtArg(num=1), attr='pt', ctx=Load())], names=[], named=[]), attr='impute', ctx=Load()), positional=[Num(n=0.0)], names=[], named=[]))
\end{lstlisting}
\end{frame}

\begin{frame}[fragile]{Working examples}
\vspace{0.3 cm}
\begin{lstlisting}[language=python]
>>> propagateTypes("""
... data.map(x => x + y)
... """,
... data=collection(integer), y=integer)
collection(integer)
\end{lstlisting}

\begin{lstlisting}[language=python]
>>> propagateTypes("""
... data.map(x => x + y)
... """,
... data=collection(integer), y=real)
collection(real)
\end{lstlisting}

\begin{lstlisting}[language=python]
>>> propagateTypes("""
... data.map(x => x + y)
... """,
... data=collection(integer, fewest=10, most=10),
... y=integer)
collection(integer, fewest=10, most=10)
\end{lstlisting}
\end{frame}

\begin{frame}[fragile]{Working examples}
\vspace{0.3 cm}
\textcolor{darkblue}{Propagate intervals of validity:}
\begin{lstlisting}[language=python]
>>> propagateTypes("""
... data.map(x => x + y)
... """,
... data=collection(real(min=3, max=5)),
... y=real(min=100, max=200))
collection(real(min=103.0, max=205.0))
\end{lstlisting}

\vfill
\textcolor{darkblue}{Properly handle shadowed variable ``x'':}
\begin{lstlisting}[language=python]
>>> propagateTypes("""
... y = x + -100;
... data.map(x => x + y)
... """,
... data=collection(real(min=3, max=5)),
... x=real(min=100, max=200))
collection(real(min=3.0, max=105.0))
\end{lstlisting}
\end{frame}

\begin{frame}{Type system is highly granular}
\vspace{0.5 cm}
Collections have \textcolor{darkblue}{fewest} and \textcolor{darkblue}{most} size attributes, and numbers have \textcolor{darkblue}{min} and \textcolor{darkblue}{max} intervals:
\begin{itemize}
\item {\tt real(almost(0), 10)} means $\{x | x \in \mathbb{R}\mbox{ and } 0 < x \le 10\}$
\item {\tt integer(almost(-inf), almost(inf))} means $\mathbb{Z}$
\item {\tt extended(-inf, inf)} means $\mathbb{R} \cup \{-\infty, \infty\}$
\item {\tt union(integer, real(0))} means

\hfill $\mathbb{Z} \cup \{x | x \in \mathbb{R}\mbox{ and } x \ge 0\}$
\end{itemize}

\vfill
\begin{uncoverenv}<2->
\textcolor{darkorange}{\bf Why?}

\vspace{0.2 cm}
To eliminate the possibility of runtime errors.

\vspace{0.2 cm}
With enough information in the type system, the compiler can identify runtime errors before submitting the job, saving the author time and protecting shared resources from waste.

\vspace{0.2 cm}
\textcolor{gray}{(Someday, they might be cloud-based and cost real money.)}
\end{uncoverenv}
\end{frame}

\begin{frame}[fragile]{Working examples}
\vspace{0.2 cm}
\begin{lstlisting}[language=python]
>>> propagateTypes("x / y", x=real, y=real)
\end{lstlisting}
\color{red}
\begin{lstlisting}[basicstyle=\ttfamily\scriptsize]
femtocode.parser.FemtocodeError: Function "/" does not accept arguments with the given types:

    /(real,
      real)

    Indeterminate form (0 / 0) is possible; constrain with if-else.

Check line:col 1:0 (pos 0):

    x / y
----^
\end{lstlisting}

\textcolor{darkblue}{Applying a constraint changes the type of ``y'' within the ``if'' and ``else'' clauses.}
\color{black}
\begin{lstlisting}[language=python]
>>> propagateTypes(
...     """if y != 0: x / y else: None""",
...     x=real, y=real)
union(null, real)
\end{lstlisting}
\end{frame}

\begin{frame}[fragile]{Working examples}
\vspace{0.5 cm}
\begin{onlyenv}<1>
\textcolor{darkblue}{``if'', ``and'', ``or'', and ``not'' propagate constraints.}
\begin{lstlisting}[language=python]
>>> propagateTypes("x == 5 and y == 6 and x == y",
... x=real, y=real)
\end{lstlisting}
\color{red}
\begin{lstlisting}[basicstyle=\ttfamily\scriptsize]
femtocode.parser.FemtocodeError: Function "==" does not accept arguments with the given types:

    ==(integer(min=5, max=5),
       integer(min=6, max=6))

    The argument types have no overlap (values can never be equal).

Check line:col 1:27 (pos 27):

    x == 5 and y == 6 and x == y
-------------------------------^
\end{lstlisting}
\end{onlyenv}
\begin{onlyenv}<2>
\textcolor{darkblue}{Order does not matter.}
\begin{lstlisting}[language=python]
>>> propagateTypes("x == y and x == 5 and y == 6",
... x=real, y=real)
\end{lstlisting}
\color{red}
\begin{lstlisting}[basicstyle=\ttfamily\scriptsize]
femtocode.parser.FemtocodeError: Function "==" does not accept arguments with the given types:

    ==(integer(min=5, max=5),
       integer(min=6, max=6))

    The argument types have no overlap (values can never be equal).

Check line:col 1:5 (pos 5):

    x == y and x == 5 and y == 6
---------^
\end{lstlisting}
\end{onlyenv}

\uncover<2>{\textcolor{gray}{(It's a short step from here to simplifying the algebra with SymPy.)}}
\end{frame}

\begin{frame}[fragile]{Working examples}
\begin{onlyenv}<1>
\begin{lstlisting}[language=python]
>>> viewAsTree("""
... a = x + y;
... b = a + y + z;
... xs.map(x => x + a + a + b).map(y => y + 2)""",
... xs=collection(real), x=real, y=real, z=real)
\end{lstlisting}
\end{onlyenv}
\begin{onlyenv}<2>
\begin{lstlisting}[language=python]
>>> viewAsStatements("""
... a = x + y;
... b = a + y + z;
... xs.map(x => x + a + a + b).map(y => y + 2)""",
... xs=collection(real), x=real, y=real, z=real)
\end{lstlisting}
\end{onlyenv}

\vspace{-0.5 cm}
\begin{minipage}[t][12cm][t]{\textwidth}
\begin{onlyenv}<1>
\begin{lstlisting}[basicstyle=\ttfamily\tiny]
Call BuiltinFunction[".map"] has type collection(real)
    Call BuiltinFunction[".map"] has type collection(real)
        Ref xs (frame None) has type collection(real)
        UserFunction has type real
            Call BuiltinFunction["+"] has type real
                Call BuiltinFunction["+"] has type real
                    Call BuiltinFunction["+"] has type real
                        Ref x (frame 2) has type real
                        Call BuiltinFunction["+"] has type real
                            Ref x (frame None) has type real
                            Ref y (frame None) has type real
                    Call BuiltinFunction["+"] has type real
                        Ref x (frame None) has type real
                        Ref y (frame None) has type real
                Call BuiltinFunction["+"] has type real
                    Call BuiltinFunction["+"] has type real
                        Call BuiltinFunction["+"] has type real
                            Ref x (frame None) has type real
                            Ref y (frame None) has type real
                        Ref y (frame None) has type real
                    Ref z (frame None) has type real
    UserFunction has type real
        Call BuiltinFunction["+"] has type real
            Ref y (frame 3) has type real
            Literal 2 has type integer(min=2, max=2)
\end{lstlisting}
\end{onlyenv}
\begin{onlyenv}<2>
\begin{lstlisting}[basicstyle=\ttfamily]
tmp_0 := (+ x y)
tmp_1 := (+ xs tmp_0)
tmp_2 := (+ x y)
tmp_3 := (+ tmp_1 tmp_2)
tmp_4 := (+ x y)
tmp_5 := (+ tmp_4 y)
tmp_6 := (+ tmp_5 z)
tmp_7 := (+ tmp_3 tmp_6)
tmp_8 := (+ tmp_7 2)
\end{lstlisting}
\end{onlyenv}
\end{minipage}
\end{frame}


\end{document}
